% !TeX program = pdflatex
\documentclass[11pt,a4paper]{article}

\usepackage[margin=1in]{geometry}
\usepackage[T1]{fontenc}
\usepackage[utf8]{inputenc}
\usepackage{lmodern}
\usepackage{microtype}
\usepackage{amsmath,amssymb,amsthm,mathtools,bm}
\usepackage{algorithm}
\usepackage{algorithmic}
\usepackage{booktabs}
\usepackage{multirow}
\usepackage{enumitem}
\usepackage{hyperref}
\usepackage[nameinlink,noabbrev]{cleveref}

\hypersetup{colorlinks=true,linkcolor=blue,citecolor=blue,urlcolor=blue}

\newtheorem{definition}{Definition}
\newtheorem{assumption}{Assumption}
\newtheorem{theorem}{Theorem}
\newtheorem{lemma}{Lemma}
\newtheorem{proposition}{Proposition}
\newtheorem{corollary}{Corollary}
\newtheorem{remark}{Remark}

\DeclareMathOperator*{\argmin}{arg\,min}
\DeclareMathOperator*{\argmax}{arg\,max}
\DeclareMathOperator{\KL}{KL}
\DeclareMathOperator{\OT}{OT}
\DeclareMathOperator{\diag}{diag}
\DeclareMathOperator{\Tr}{Tr}
\DeclareMathOperator{\Lip}{Lip}
\DeclareMathOperator{\conv}{conv}
\newcommand{\E}{\mathbb{E}}
\newcommand{\Pbb}{\mathbb{P}}
\newcommand{\R}{\mathbb{R}}
\newcommand{\N}{\mathbb{N}}
\newcommand{\cY}{\mathcal{Y}}
\newcommand{\cX}{\mathcal{X}}
\newcommand{\cP}{\mathcal{P}}
\newcommand{\cH}{\mathcal{H}}
\newcommand{\cD}{\mathcal{D}}
\newcommand{\cL}{\mathcal{L}}
\newcommand{\cU}{\mathcal{U}}
\newcommand{\cB}{\mathcal{B}}
\newcommand{\what}{\widehat}
\newcommand{\wt}{\widetilde}
\newcommand{\bs}{\boldsymbol}
\newcommand{\norm}[1]{\left\lVert #1\right\rVert}
\newcommand{\abs}[1]{\left\lvert #1\right\rvert}
\newcommand{\ip}[2]{\left\langle #1,#2\right\rangle}

\title{Sinkhorn-Minimax Active Learning on Distributions: \\
A Geometric Distributionally Robust Reformulation of Active Robust Learning}
\author{Seyed Hossein Ghafarian and Hadi Sadoghi Yazdi\\
\small Department of Computer Engineering, Ferdowsi University of Mashhad, Mashhad, Iran}
\date{}

\begin{document}
\maketitle

\begin{abstract}
Learning on distributions considers training examples that are not single vectors but finite samples drawn from latent probability measures. Active learning in this setting is difficult because querying a label for a distributional example must account for two coupled uncertainties: label uncertainty and distribution-estimation uncertainty. Existing active robust learning on distributions addresses this issue through kernel mean embeddings and Bayesian uncertainty in a reproducing kernel Hilbert space. This paper develops a different and minimally overlapping perspective: active learning on distributions is formulated as a geometric distributionally robust optimization problem over ambiguity balls defined by entropic optimal transport. The resulting method, called Sinkhorn-Minimax Active Learning on Distributions (S-MALD), queries the distributional example that minimizes a worst-case post-query risk upper bound under Sinkhorn/Wasserstein perturbations of the empirical distributional population. We derive a post-query risk bound, a finite-sample ambiguity radius principle for bags of unequal size, a reduction to a tractable quadratic label objective through a Sinkhorn affinity matrix, and a limiting connection showing how the older kernel-embedding view appears as a computational surrogate of a geometric transport problem. The proposed formulation preserves the minimax spirit of probabilistic minimax active learning while replacing the RKHS-only representation with a geometry-aware ambiguity model that is more suitable for modern optimal-transport-based learning.
\end{abstract}

\noindent\textbf{Keywords:} active learning; learning on distributions; Sinkhorn divergence; Wasserstein distributionally robust optimization; optimal transport; minimax learning; distributional data.

\section{Introduction}

Many modern data sets are naturally organized as \emph{sets}, \emph{bags}, or \emph{empirical distributions}. Examples include a document represented by its sentence embeddings, an image object represented by local patches, a user represented by a set of interactions, a patient represented by repeated measurements, or a sensor region represented by a finite point cloud. In these problems the training example is not a vector $x\in\R^d$ but a probability distribution $P\in\cP(\cX)$, observed only through a finite sample
\[
X_i=\{x_{i1},\ldots,x_{in_i}\}\sim P_i^{n_i}.
\]
The learning task is to infer a classifier $g:\cP(\cX)\to\cY$ from labeled distributional examples and a pool of unlabeled distributional examples.

Active learning is especially attractive in this setting because the label of a distributional example may be expensive. However, the query decision is more delicate than in ordinary vector-based active learning. A distributional example can be uncertain because its class is ambiguous, but it can also be uncertain because $P_i$ is poorly estimated from a small or noisy bag $X_i$. A query rule that ignores the second source of uncertainty may over-query unreliable bags or under-query distributional examples that are geometrically informative.

Earlier active robust learning on distributions formulated this difficulty through kernel mean embeddings (KMEs). A probability measure $P$ is mapped to an RKHS element $\mu_P=\E_{X\sim P}k(X,\cdot)$, and uncertainty is handled through empirical or Bayesian estimates of this embedding. This is elegant and computationally useful, but it makes the Hilbert representation the central object of the theory. In contrast, the present article asks whether the same active-learning principle can be expressed directly in the geometry of probability measures.

The answer is affirmative. We formulate active learning on distributions as a distributionally robust optimization (DRO) problem over ambiguity sets centered at empirical bags. Instead of measuring distributional perturbations only by $\norm{\mu_P-\mu_Q}_{\cH}$, we use entropic optimal transport and its debiased Sinkhorn divergence. This gives a geometry-aware active learning rule: the selected query is the one whose label most reduces the worst-case post-query risk over plausible perturbations of the empirical distributional population.

\paragraph{Main idea.}
Let $L$ and $U$ denote the labeled and unlabeled distributional examples, and let $Q\subseteq U$ be a candidate query batch. The ideal query should solve a problem of the form
\begin{equation}
Q^*\in \argmin_{Q\subseteq U,\,|Q|=b}\;
\max_{\bm y_Q\in\cY^b}\;
\sup_{M\in\cB_\varepsilon(\widehat M,\rho)}
\inf_{\bm y_{U\setminus Q}\in\cY^{|U|-b}}
\mathcal{R}(g_{L\cup Q};M,\bm y_Q,\bm y_{U\setminus Q}),
\label{eq:ideal-sinkhorn-pmal}
\end{equation}
where $\widehat M$ is the empirical population of observed distributional examples, $\cB_\varepsilon(\widehat M,\rho)$ is a Sinkhorn/entropic-Wasserstein ambiguity set, and $\mathcal{R}$ is a post-query risk surrogate. The minimization over $Q$, maximization over unknown query labels, supremum over plausible distributional perturbations, and infimum over remaining unlabeled labels together yield a robust geometric version of probabilistic minimax active learning.

\paragraph{Contributions.}
This paper makes the following contributions.
\begin{enumerate}[leftmargin=1.4em]
    \item We introduce a Sinkhorn/Wasserstein DRO formulation for active learning on distributions, where ambiguity is placed on distributions of distributional examples rather than on vector inputs.
    \item We derive a post-query risk upper bound showing how the classical PMAL criterion acquires an explicit transport-robustness penalty.
    \item We propose S-MALD, a tractable Sinkhorn-minimax query rule based on entropic transport between bags and a quadratic label objective.
    \item We provide a finite-sample radius principle for unequal bag sizes, making the query score sensitive to distributional examples with small $n_i$.
    \item We clarify the relation to earlier KME-based active robust learning: KME/MMD appears as a limiting or surrogate representation, whereas the new formulation is driven by geometric ambiguity in probability space.
\end{enumerate}

\section{Relation to the Previous Cybernetics Paper}

The intended novelty of this manuscript is not to rewrite the earlier active robust learning paper with different notation. The difference is structural. The previous work uses PMAL and PMARL to handle inexact examples, then specializes the framework to learning on distributions through Bayesian kernel embedding. The present work keeps the minimax active-learning philosophy but changes the mathematical object on which robustness is defined: the ambiguity set is now a Sinkhorn/Wasserstein ball over empirical measures of bags.

\begin{table}[t]
\centering
\caption{Non-overlap between the earlier KME-based article and the proposed Sinkhorn-DRO article.}
\label{tab:novelty}
\begin{tabular}{p{0.28\linewidth}p{0.31\linewidth}p{0.31\linewidth}}
\toprule
Aspect & Earlier KME/PMARLDB direction & Proposed S-MALD direction \\
\midrule
Representation of a distributional example & RKHS mean embedding $\mu_P$ and Bayesian uncertainty of $\mu_P$ & Empirical measure $\widehat P_i$ compared through entropic optimal transport \\
Uncertainty model & Covariance of the estimated kernel embedding & Sinkhorn/Wasserstein ambiguity ball around empirical bags \\
Main geometry & Hilbert-space norm and kernel similarity & Transport geometry on probability measures \\
Query criterion & PMAL/PMARL quadratic objective with embedding uncertainty & Sinkhorn-minimax post-query risk bound with transport robustness penalty \\
Computational object & Kernel matrix from expected kernels or Bayesian KME & Sinkhorn cost/divergence matrix and Sinkhorn affinity \\
Scientific claim & Robust active learning under inexact distribution embeddings & Distributionally robust active learning under geometric perturbations of distributions \\
\bottomrule
\end{tabular}
\end{table}

\section{Problem Setup}

Let $\cX\subset\R^d$ be a compact instance space and $\cY=\{1,\ldots,C\}$ a finite label set. A distributional example is a latent probability measure $P_i\in\cP(\cX)$ observed through a finite sample
\[
X_i=\{x_{i1},\ldots,x_{in_i}\},\qquad x_{ij}\sim P_i.
\]
The empirical distribution associated with the $i$th bag is
\begin{equation}
\widehat P_i = \sum_{j=1}^{n_i} a_{ij}\delta_{x_{ij}},
\qquad a_{ij}\ge 0,
\qquad \sum_{j=1}^{n_i}a_{ij}=1.
\label{eq:empirical-bag}
\end{equation}
Uniform weights $a_{ij}=1/n_i$ are used unless bag instances have reliability weights.

At an active learning round, the index set is divided into labeled and unlabeled parts, $L$ and $U$. Labels $\bm y_L$ are known and labels $\bm y_U$ are unknown. The active learner selects a query batch $Q\subseteq U$ of size $b$ and asks an oracle for $\bm y_Q$. The objective is to select $Q$ such that the classifier trained after receiving $\bm y_Q$ has small risk.

Let $\eta^*(P,c)=\Pbb(Y=c\mid P)$ be the true posterior probability and $\eta_D(P,c)$ the posterior estimate trained from labeled data $D$. The plug-in classifier is
\begin{equation}
g_D(P)=\argmax_{c\in\cY}\eta_D(P,c).
\end{equation}
The excess plug-in error at $P$ is controlled by the posterior approximation error
\begin{equation}
A_D(P)=\max_{c\in\cY}\{\eta^*(P,c)-\eta_D(P,c)\}_+.
\label{eq:posterior-error}
\end{equation}
For a population measure $M$ over distributional examples, define
\begin{equation}
\Delta_D(M)=\int_{\cP(\cX)}A_D(P)\,dM(P).
\label{eq:delta-D}
\end{equation}
The purpose of active learning is to query examples that reduce a bound on $\Delta_D(M)$ before their labels are known.

\section{Entropic Optimal Transport Between Distributional Examples}

Let $c:\cX\times\cX\to\R_+$ be a ground cost, for example $c(x,z)=\norm{x-z}_2^2$ or $c_\phi(x,z)=\norm{\phi(x)-\phi(z)}_2^2$ for a feature extractor $\phi$. For two empirical bags $\widehat P_i=\sum_u a_{iu}\delta_{x_{iu}}$ and $\widehat P_j=\sum_v a_{jv}\delta_{x_{jv}}$, let $C_{ij}\in\R^{n_i\times n_j}$ be the cost matrix with entries $C_{ij}(u,v)=c(x_{iu},x_{jv})$.

\begin{definition}[Entropic optimal transport]
For $\varepsilon>0$, the entropic OT cost between two empirical bags is
\begin{equation}
\OT_\varepsilon(\widehat P_i,\widehat P_j)
=\\min_{\Gamma\in\Pi(a_i,a_j)}
\ip{\Gamma}{C_{ij}}+\varepsilon\KL(\Gamma\,\|\,a_ia_j^\top),
\label{eq:entropic-ot}
\end{equation}
where $\Pi(a_i,a_j)=\{\Gamma\ge 0: \Gamma\mathbf{1}=a_i,\Gamma^\top\mathbf{1}=a_j\}$.
\end{definition}

\begin{definition}[Sinkhorn divergence]
The debiased Sinkhorn divergence is
\begin{equation}
S_\varepsilon(\widehat P_i,\widehat P_j)
=\OT_\varepsilon(\widehat P_i,\widehat P_j)
-\frac{1}{2}\OT_\varepsilon(\widehat P_i,\widehat P_i)
-\frac{1}{2}\OT_\varepsilon(\widehat P_j,\widehat P_j).
\label{eq:sinkhorn-divergence}
\end{equation}
\end{definition}

The role of $\varepsilon$ is central. Small $\varepsilon$ behaves closer to classical optimal transport and preserves geometric displacement. Larger $\varepsilon$ smooths the problem, improves numerical stability, and creates a bridge toward kernel/MMD-like behavior. Thus $\varepsilon$ is not merely a numerical parameter; it controls the geometry of the active learning criterion.

\section{Distributional Ambiguity at Two Levels}

In learning on distributions there are two empirical approximations:
\begin{enumerate}[leftmargin=1.4em]
    \item each latent example $P_i$ is approximated by a finite bag $\widehat P_i$;
    \item the population of distributional examples is approximated by a finite pool $\widehat M=\sum_{i=1}^N \omega_i\delta_{\widehat P_i}$.
\end{enumerate}
The first approximation concerns within-bag uncertainty; the second concerns population uncertainty. We encode both in a transport ambiguity set.

Let $d_\varepsilon(P,Q)$ denote either $\OT_\varepsilon(P,Q)^{1/2}$ or the debiased $S_\varepsilon(P,Q)^{1/2}$ whenever the latter is nonnegative. Define a cost between distributional examples by
\begin{equation}
\bar c_\varepsilon(P,Q)=d_\varepsilon(P,Q)^2.
\end{equation}
For two probability measures $M,N\in\cP(\cP(\cX))$ over distributional examples, define a second-level transport discrepancy
\begin{equation}
\mathcal{W}_{\varepsilon}(M,N)
=\inf_{\Pi\in\Pi(M,N)}\int \bar c_\varepsilon(P,Q)\,d\Pi(P,Q).
\label{eq:second-level-ot}
\end{equation}
The ambiguity ball centered at the empirical pool is
\begin{equation}
\cB_\varepsilon(\widehat M,\rho)=
\{M\in\cP(\cP(\cX)): \mathcal{W}_{\varepsilon}(M,\widehat M)\le \rho\}.
\label{eq:ambiguity-ball}
\end{equation}
This formulation places uncertainty directly on the population of distributions and does not require the distributions to be embedded first into an RKHS.

\section{Sinkhorn-Minimax Active Learning}

\subsection{Ideal Robust Query Objective}

For a candidate query batch $Q\subseteq U$, the post-query data set is $L_Q=L\cup Q$, but $\bm y_Q$ is not yet known. Following the minimax active learning principle, we consider the worst possible query labels and the most favorable completion of the remaining unlabeled labels in the risk bound. The robust post-query objective is
\begin{equation}
\mathfrak{A}_\varepsilon(Q)
=\max_{\bm y_Q\in\cY^{|Q|}}
\sup_{M\in\cB_\varepsilon(\widehat M,\rho_Q)}
\min_{\bm y_{U\setminus Q}\in\cY^{|U|-|Q|}}
\Delta_{L_Q}(M).
\label{eq:robust-acquisition-general}
\end{equation}
The query rule is
\begin{equation}
Q^*\in\argmin_{Q\subseteq U,\,|Q|=b}\mathfrak{A}_\varepsilon(Q).
\label{eq:s-mald-query}
\end{equation}
This is the geometric counterpart of probabilistic minimax active learning. The difference is that the uncertainty is not an additive vector uncertainty or a covariance around an RKHS embedding; it is a distributional ambiguity set induced by transport geometry.

\subsection{Risk Bound}

\begin{assumption}[Transport Lipschitz posterior error]
For any labeled set $D$, the posterior error $A_D(P)$ in \eqref{eq:posterior-error} is $L_D$-Lipschitz with respect to $d_\varepsilon$, namely
\begin{equation}
\abs{A_D(P)-A_D(Q)}\le L_D d_\varepsilon(P,Q),
\qquad \forall P,Q\in\cP(\cX).
\label{eq:lipschitz-error}
\end{equation}
\end{assumption}

\begin{theorem}[Transport-robust post-query risk bound]
\label{thm:robust-risk-bound}
Assume \cref{eq:lipschitz-error}. For any candidate query batch $Q$ and any ambiguity radius $\rho_Q$, every $M\in\cB_\varepsilon(\widehat M,\rho_Q)$ satisfies
\begin{equation}
\Delta_{L_Q}(M)
\le \Delta_{L_Q}(\widehat M)+L_{L_Q}\rho_Q^{1/2}.
\label{eq:delta-transport-bound}
\end{equation}
Consequently, if the plug-in excess risk is bounded by $L_D-L^*\le \zeta\Delta_D(M)$ for a constant $\zeta\ge 1$, then
\begin{equation}
L_{L_Q}-L^*
\le \zeta\left[\Delta_{L_Q}(\widehat M)+L_{L_Q}\rho_Q^{1/2}\right].
\label{eq:excess-bound}
\end{equation}
\end{theorem}

\begin{proof}
Let $\Pi$ be any coupling between $M$ and $\widehat M$. By \cref{eq:lipschitz-error},
\[
\int A_{L_Q}(P)dM(P)-\int A_{L_Q}(Q')d\widehat M(Q')
\le L_{L_Q}\int d_\varepsilon(P,Q')d\Pi(P,Q').
\]
Applying Cauchy--Schwarz and taking the infimum over couplings gives
\[
\Delta_{L_Q}(M)-\Delta_{L_Q}(\widehat M)
\le L_{L_Q}\mathcal{W}_\varepsilon(M,\widehat M)^{1/2}
\le L_{L_Q}\rho_Q^{1/2}.
\]
Combining this inequality with the standard plug-in excess-risk bound proves the claim.
\end{proof}

\begin{corollary}[Computable robust acquisition bound]
\label{cor:computable-acquisition}
Under the assumptions of \cref{thm:robust-risk-bound}, it is sufficient to minimize
\begin{equation}
\overline{\mathfrak{A}}_\varepsilon(Q)
=\max_{\bm y_Q}
\min_{\bm y_{U\setminus Q}}
\Delta_{L_Q}(\widehat M)
+ L_{L_Q}\rho_Q^{1/2}.
\label{eq:computable-bound}
\end{equation}
The first term is the ordinary post-query minimax term, while the second term is the explicit transport-robustness price.
\end{corollary}

\subsection{Finite-Sample Radius for Unequal Bags}

The radius $\rho_Q$ should increase when candidate queried distributions are poorly estimated. This is important because distributional active learning often contains bags with unequal sample sizes.

\begin{assumption}[Bag-level concentration]
For each $i$ and confidence level $\delta_i\in(0,1)$, there exists a nonincreasing function $r(n_i,\delta_i)$ such that
\begin{equation}
d_\varepsilon(P_i,\widehat P_i)\le r(n_i,\delta_i)
\label{eq:bag-concentration}
\end{equation}
with probability at least $1-\delta_i$.
\end{assumption}

\begin{proposition}[Pool ambiguity radius]
\label{prop:radius}
Let $\widehat M=\sum_{i=1}^N\omega_i\delta_{\widehat P_i}$ and $M^*=\sum_{i=1}^N\omega_i\delta_{P_i}$. Under the bag-level concentration assumption, with probability at least $1-\sum_i\delta_i$,
\begin{equation}
\mathcal{W}_\varepsilon(M^*,\widehat M)
\le \sum_{i=1}^N \omega_i r(n_i,\delta_i)^2.
\label{eq:pool-radius}
\end{equation}
Thus a natural data-dependent radius is
\begin{equation}
\rho=\sum_{i=1}^N\omega_i r(n_i,\delta_i)^2.
\label{eq:data-dependent-radius}
\end{equation}
\end{proposition}

\begin{proof}
Use the coupling that transports each atom $P_i$ of $M^*$ to its empirical counterpart $\widehat P_i$ with mass $\omega_i$. The transport cost under this coupling is $\sum_i\omega_i d_\varepsilon(P_i,\widehat P_i)^2$, and the stated inequality follows from the concentration assumption and the union bound.
\end{proof}

In practice, $r(n_i,\delta_i)$ can be set from a theoretical concentration inequality or estimated by within-bag bootstrap. This makes the algorithm conservative for small or heterogeneous bags without discarding them.

\section{Tractable Sinkhorn-Minimax Approximation}

The exact objective \eqref{eq:robust-acquisition-general} is intractable for large pools. We derive a practical approximation that mirrors the quadratic label form of PMAL while replacing the kernel matrix by a Sinkhorn-induced affinity matrix.

\subsection{Sinkhorn Affinity Matrix}

For all pairs of distributional examples, compute
\begin{equation}
D^\varepsilon_{ij}=S_\varepsilon(\widehat P_i,\widehat P_j).
\end{equation}
A Sinkhorn affinity matrix is then defined by
\begin{equation}
K^\varepsilon_{ij}=\exp(-\gamma D^\varepsilon_{ij}),
\label{eq:sinkhorn-kernel}
\end{equation}
where $\gamma>0$ is a scale parameter. If $K^\varepsilon$ is not positive semidefinite in finite precision, we use the projected matrix
\begin{equation}
K^\varepsilon_+ = V\max(\Lambda,0)V^\top,
\qquad K^\varepsilon=V\Lambda V^\top.
\label{eq:psd-projection}
\end{equation}
This projection is only a computational safeguard. The conceptual object remains the Sinkhorn geometry between bags.

\subsection{Quadratic Label Surrogate}

Let $K=K^\varepsilon_+$ and define the regularized smoothing matrix
\begin{equation}
H_\varepsilon=(K+\lambda I)^{-1}K(K+\lambda I)^{-1}.
\label{eq:H-eps}
\end{equation}
For a partition into labeled indices $L$, candidate query indices $Q$, and remaining unlabeled indices $R=U\setminus Q$, let $\bm y=(\bm y_L,\bm y_Q,\bm y_R)$. A quadratic approximation to the negative log-posterior label likelihood is
\begin{equation}
\mathcal{J}_\varepsilon(\bm y;Q)=\bm y^\top S_\varepsilon^{(Q)}\bm y,
\label{eq:quad-label-objective}
\end{equation}
where $S_\varepsilon^{(Q)}$ is a positive semidefinite matrix obtained from $H_\varepsilon$ and the bag uncertainty weights. A simple choice is
\begin{equation}
S_\varepsilon^{(Q)}=\\Upsilon_Q-
\Upsilon_Q\left(\lambda^{-1}I+K\Upsilon_Q\right)^{-1}K\Upsilon_Q,
\label{eq:S-eps}
\end{equation}
with
\begin{equation}
\Upsilon_Q=\diag(\kappa_i),
\qquad
\kappa_i=\left(\sigma^2+\tau_i\right)^{-1},
\qquad
\tau_i=\chi\,r(n_i,\delta_i)^2.
\label{eq:uncertainty-weight}
\end{equation}
Here $\tau_i$ penalizes distributional examples whose bags are poorly estimated, and $\chi>0$ controls the strength of robust correction.

The tractable S-MALD acquisition is
\begin{equation}
\widehat{\mathfrak{A}}_\varepsilon(Q)
=\max_{\bm y_Q\in\cY^{|Q|}}
\min_{\bm y_R\in\conv(\cY)^{|R|}}
\bm y^\top S_\varepsilon^{(Q)}\bm y
+\lambda_\rho \left(\sum_{i\in Q}\omega_i r(n_i,\delta_i)^2\right)^{1/2}.
\label{eq:tractable-smald}
\end{equation}
For binary labels, $\conv(\cY)=[-1,1]$. The inner minimization is convex when $S_\varepsilon^{(Q)}\succeq 0$. For batch size $b=1$, evaluating all candidates is practical.

\begin{algorithm}[t]
\caption{Sinkhorn-Minimax Active Learning on Distributions (S-MALD)}
\label{alg:smald}
\begin{algorithmic}[1]
\REQUIRE Labeled bags $\{(X_i,y_i)\}_{i\in L}$, unlabeled bags $\{X_i\}_{i\in U}$, batch size $b$, entropic parameter $\varepsilon$, scale $\gamma$, regularization $\lambda$.
\STATE Construct empirical measures $\widehat P_i$ for all bags using \eqref{eq:empirical-bag}.
\STATE Compute pairwise Sinkhorn divergences $D^\varepsilon_{ij}=S_\varepsilon(\widehat P_i,\widehat P_j)$.
\STATE Form the Sinkhorn affinity matrix $K^\varepsilon_{ij}=\exp(-\gamma D^\varepsilon_{ij})$ and project it to $K^\varepsilon_+$ if necessary.
\STATE Estimate bag uncertainty radii $r(n_i,\delta_i)$ by theory or bootstrap.
\STATE Build $S_\varepsilon^{(Q)}$ from \eqref{eq:S-eps} and \eqref{eq:uncertainty-weight}.
\FOR{each candidate query batch $Q\subseteq U$, $|Q|=b$}
    \STATE Compute $\widehat{\mathfrak{A}}_\varepsilon(Q)$ in \eqref{eq:tractable-smald}.
\ENDFOR
\STATE Query $Q^*=\argmin_Q \widehat{\mathfrak{A}}_\varepsilon(Q)$.
\STATE Ask the oracle for labels $\bm y_{Q^*}$ and update $L\leftarrow L\cup Q^*$, $U\leftarrow U\setminus Q^*$.
\end{algorithmic}
\end{algorithm}

\section{Connection to Kernel Mean Embedding}

The proposed method is not an arbitrary replacement of KME by Wasserstein distance. Instead, it creates a continuum between geometric optimal transport and kernel discrepancy.

\begin{proposition}[KME/MMD as a limiting surrogate]
\label{prop:kme-limit}
Let $k$ be a characteristic kernel and let $\mathrm{MMD}_k(P,Q)=\norm{\mu_P-\mu_Q}_{\cH}$ be the maximum mean discrepancy. Sinkhorn-type divergences form a regularized family of geometric discrepancies that interpolates between optimal transport behavior at small entropic regularization and kernel/MMD-like behavior at strong regularization. Therefore, a KME-based active learner can be interpreted as a Hilbertian surrogate of a transport-robust active learner, whereas S-MALD keeps the transport geometry explicit through $S_\varepsilon$.
\end{proposition}

\begin{remark}
This proposition is the conceptual reason why the current manuscript has low overlap with the earlier KME article. We do not claim that the previous KME objective was wrong. Rather, we show that it can be understood as one endpoint or surrogate of a broader Sinkhorn/Wasserstein active-learning geometry.
\end{remark}

\section{Theoretical Consequences}

\subsection{Query Stability}

\begin{theorem}[Stability of the Sinkhorn-minimax score]
\label{thm:stability}
Assume that the Sinkhorn divergence estimates satisfy
\begin{equation}
\max_{i,j}\abs{D^\varepsilon_{ij}-\widetilde D^\varepsilon_{ij}}\le \delta_D.
\end{equation}
Then, for $K^\varepsilon_{ij}=\exp(-\gamma D^\varepsilon_{ij})$,
\begin{equation}
\max_{i,j}\abs{K^\varepsilon_{ij}-\widetilde K^\varepsilon_{ij}}
\le \gamma \delta_D.
\label{eq:kernel-perturbation}
\end{equation}
Moreover, if $K^\varepsilon+\lambda I\succeq \lambda I$ and the label vectors are bounded as $\norm{\bm y}\le B_y$, then the quadratic score perturbation satisfies
\begin{equation}
\abs{\bm y^\top S_\varepsilon\bm y-\bm y^\top \widetilde S_\varepsilon\bm y}
\le C(\lambda,\gamma,B_y,N)\delta_D
\end{equation}
for a finite constant $C$.
\end{theorem}

\begin{proof}
The exponential map is $\gamma$-Lipschitz on $\R_+$ after bounding $e^{-\gamma t}\le 1$, which proves \eqref{eq:kernel-perturbation}. The matrix $S_\varepsilon$ is obtained from additions, products, and inverses of matrices whose smallest eigenvalue is bounded below by $\lambda$. Standard perturbation bounds for matrix inversion give the stated Lipschitz dependence of $S_\varepsilon$ on $K^\varepsilon$.
\end{proof}

\subsection{Consistency of the Robust Query Objective}

\begin{theorem}[Asymptotic recovery of the non-robust objective]
\label{thm:consistency}
Assume that $r(n_i,\delta_i)\to 0$ for all queried bags as $n_i\to\infty$, and that $\rho_t\to 0$ as the active learning round $t$ increases. Then the robust acquisition bound \eqref{eq:computable-bound} converges to the non-robust PMAL-style bound:
\begin{equation}
\overline{\mathfrak{A}}_\varepsilon(Q)-
\max_{\bm y_Q}\min_{\bm y_{U\setminus Q}}\Delta_{L_Q}(\widehat M)
\to 0.
\end{equation}
If the underlying plug-in learner is consistent on distributional examples, then the corresponding excess risk converges to the Bayes risk under the same conditions required by the base learner.
\end{theorem}

\begin{proof}
The difference between the robust and non-robust objectives is exactly the transport penalty $L_{L_Q}\rho_Q^{1/2}$ in \eqref{eq:computable-bound}. Under the stated assumptions this term vanishes. Consistency follows by combining this vanishing robust penalty with the consistency of the base plug-in learner.
\end{proof}

\section{Experimental Design for a Low-Overlap Submission}

The experiments should not simply reproduce the previous Cybernetics paper. The new manuscript should evaluate the specific scientific claim of this paper: geometric ambiguity improves active learning on distributions when the bags are finite, heterogeneous, and geometrically shifted.

\subsection{Recommended Data Sets}

The following experiments are suitable because they match distributional learning but test Sinkhorn geometry rather than only KME uncertainty.
\begin{enumerate}[leftmargin=1.4em]
    \item \textbf{Synthetic Gaussian distributions.} Each example is a Gaussian or Gaussian mixture represented by a finite bag. Vary mean shifts, covariance shifts, and bag size imbalance.
    \item \textbf{20 Newsgroups as distributions of sentence/word embeddings.} Each document or class-fragment is represented as a bag of embeddings; labels are topic categories.
    \item \textbf{COIL object views.} Each object is represented by a distribution of views or image descriptors.
    \item \textbf{USPS or MNIST bags.} Each distributional example is a bag of local descriptors or augmentations.
\end{enumerate}

\subsection{Baselines}

To make the novelty clear, the comparison should include both classical active learning and distributional baselines:
\begin{enumerate}[leftmargin=1.4em]
    \item random query selection;
    \item uncertainty sampling;
    \item margin-based active learning;
    \item KME/MMD-based active learning;
    \item PMARLDS/PMARLDB-style robust active learning if available;
    \item Wasserstein active learning without entropic regularization;
    \item proposed S-MALD with an $\varepsilon$ sweep.
\end{enumerate}

\subsection{Ablation Studies}

The paper should include the following ablations.
\begin{enumerate}[leftmargin=1.4em]
    \item \textbf{Effect of $\varepsilon$.} Show the transition from transport-like behavior to MMD-like behavior.
    \item \textbf{Effect of bag size.} Compare equal-size and unequal-size bags; S-MALD should benefit from the radius $r(n_i,\delta_i)$.
    \item \textbf{Effect of geometric shift.} Perturb bags by translation, rotation, or covariance changes and compare against KME-only criteria.
    \item \textbf{Runtime.} Report the cost of Sinkhorn computation and the benefit of entropic regularization.
\end{enumerate}

\subsection{Suggested Reporting Template}

Use accuracy versus number of queried labels as the main plot. Report the area under the active learning curve (AUALC), final accuracy, and mean rank across data sets. For robustness, report performance under bag-size imbalance and distributional shift. This will make the paper distinct from the previous article, whose central empirical claim was the benefit of Bayesian kernel embedding.

\section{Discussion}

The main scientific advantage of S-MALD is that it moves the uncertainty model from a Hilbert-space representation to a geometric ambiguity set over distributions. This is important for applications where the geometry of the instance space matters: document embeddings, images, point clouds, sensor distributions, and any data where moving mass from one region to another is more meaningful than comparing only feature averages.

The method also preserves the attractive minimax structure of PMAL. Before a label is queried, its value is unknown, so a worst-case label view is appropriate. However, the older formulation only partially captured the fact that distributional examples are themselves uncertain. S-MALD makes this explicit through a robust transport penalty that increases for small or unreliable bags.

\section{Conclusion}

This paper proposed Sinkhorn-Minimax Active Learning on Distributions, a geometric distributionally robust reformulation of active robust learning. Instead of treating kernel mean embedding as the primary representation of distributions, the proposed method defines uncertainty through entropic optimal transport ambiguity sets and derives a robust post-query risk bound. A practical quadratic approximation was introduced through a Sinkhorn affinity matrix and a finite-sample bag uncertainty radius. The framework explains KME/MMD-based active learning as a limiting or surrogate case while opening a lower-overlap research direction based on modern optimal transport, Sinkhorn divergence, and distributionally robust active learning.

\section*{Acknowledgment}
The authors may add funding, institutional, or thesis-related acknowledgments here.

\begin{thebibliography}{99}

\bibitem{settles2012active}
B. Settles, \emph{Active Learning}. Morgan \& Claypool, 2012.

\bibitem{muandet2012smm}
K. Muandet, K. Fukumizu, F. Dinuzzo, and B. Schölkopf, ``Learning from distributions via support measure machines,'' in \emph{Advances in Neural Information Processing Systems}, 2012.

\bibitem{muandet2017kme}
K. Muandet, K. Fukumizu, B. Sriperumbudur, and B. Schölkopf, ``Kernel mean embedding of distributions: A review and beyond,'' \emph{Foundations and Trends in Machine Learning}, vol. 10, no. 1--2, pp. 1--141, 2017.

\bibitem{szabo2016distribution}
Z. Szabó, B. Sriperumbudur, B. Póczos, and A. Gretton, ``Learning theory for distribution regression,'' \emph{Journal of Machine Learning Research}, vol. 17, no. 152, pp. 1--40, 2016.

\bibitem{flaxman2016bayesian}
S. Flaxman, Y. W. Teh, and D. Sejdinovic, ``Bayesian learning of kernel embeddings,'' in \emph{Proceedings of the Thirty-Second Conference on Uncertainty in Artificial Intelligence}, 2016.

\bibitem{ghafarian_cybernetics}
S. H. Ghafarian and H. Sadoghi Yazdi, ``Prepare for the worst, hope for the best: Active robust learning on distributions,'' manuscript/published article in cybernetics venue.

\bibitem{cuturi2013sinkhorn}
M. Cuturi, ``Sinkhorn distances: Lightspeed computation of optimal transport,'' in \emph{Advances in Neural Information Processing Systems}, 2013.

\bibitem{peyre2019computational}
G. Peyré and M. Cuturi, ``Computational optimal transport: With applications to data science,'' \emph{Foundations and Trends in Machine Learning}, vol. 11, no. 5--6, pp. 355--607, 2019.

\bibitem{feydy2019sinkhorn}
J. Feydy, T. Séjourné, F.-X. Vialard, S. Amari, A. Trouvé, and G. Peyré, ``Interpolating between optimal transport and MMD using Sinkhorn divergences,'' in \emph{Proceedings of AISTATS}, 2019.

\bibitem{kuhn2019wasserstein}
D. Kuhn, P. Mohajerin Esfahani, V. A. Nguyen, and S. Shafieezadeh-Abadeh, ``Wasserstein distributionally robust optimization: Theory and applications in machine learning,'' 2019.

\bibitem{wang2021sinkhornDRO}
J. Wang, R. Gao, and Y. Xie, ``Sinkhorn distributionally robust optimization,'' arXiv:2109.11926, 2021.

\bibitem{li2021hilbert}
Q. Li, Z. Wang, G. Li, and J. Pang, ``Hilbert Sinkhorn divergence for optimal transport,'' in \emph{Proceedings of CVPR}, 2021.

\bibitem{takeno2025dral}
S. Takeno et al., ``Distributionally robust active learning for Gaussian process regression,'' arXiv:2502.16870, 2025.

\end{thebibliography}

\end{document}
